\documentclass[a4paper,12pt]{article}
\usepackage[T2A]{fontenc}
\usepackage[utf8]{inputenc}
\usepackage[russian]{babel}
\usepackage{amsmath,amssymb,mathtools}
\usepackage{helvet}
\renewcommand{\familydefault}{\sfdefault}
\usepackage{geometry}
\geometry{left=23mm,right=23mm,top=22mm,bottom=23mm}
\usepackage{microtype}
\usepackage{icomma}
\usepackage{xcolor}
\usepackage{booktabs,tabularx,array}
\usepackage{enumitem}
\usepackage{fancyhdr}
\usepackage{tikz}
\usetikzlibrary{arrows.meta,positioning,shapes.geometric,calc}
\usepackage{setspace}
\usepackage{hyperref}
\hypersetup{hidelinks}
\definecolor{accent}{HTML}{008080}
\definecolor{darkaccent}{HTML}{005B5B}
\definecolor{textgray}{HTML}{333333}
\definecolor{lightgray}{HTML}{728080}
\color{textgray}
\onehalfspacing
\setlength{\parindent}{0pt}
\setlength{\parskip}{4.5pt}
\setcounter{secnumdepth}{1}
\setlist[enumerate]{leftmargin=7mm,itemsep=4pt,topsep=4pt}
\pagestyle{fancy}
\fancyhf{}
\fancyhead[L]{\small\color{lightgray}Прикладные задачи и встречное движение}
\fancyhead[R]{\small\color{lightgray}Анна Баннова}
\fancyfoot[C]{\small\thepage}
\renewcommand{\headrulewidth}{0.35pt}
\renewcommand{\headrule}{\hbox to\headwidth{\color{accent}\leaders\hrule height \headrulewidth\hfill}}
\renewcommand{\footrulewidth}{0pt}
\setlength{\headheight}{25pt}
\addtolength{\topmargin}{-10pt}
\newcommand{\key}[1]{\textcolor{darkaccent}{#1}}
\newcommand{\lessonsection}[1]{\section{\texorpdfstring{#1}{#1}}\vspace{-3pt}\noindent\color{accent}\rule{\linewidth}{0.4pt}\color{textgray}\vspace{2pt}}
\newcommand{\example}[1]{\par\medskip\textcolor{darkaccent}{\textbf{#1}}\par\smallskip}
\newcommand{\task}[2]{\par\smallskip\noindent\textbf{#1.}\quad #2\par}
\newcommand{\answerrow}[2]{#1 & #2\\}
\begin{document}
\thispagestyle{empty}
\begin{center}
\vspace*{13mm}
{\Large\color{darkaccent}МАТЕМАТИКА\par}
\vspace{17mm}
{\huge Чек-лист\par}
\vspace{12mm}
{\LARGE\color{accent}Прикладные формулы\\[4pt] и текстовые задачи\\[4pt] на встречное движение\par}
\vspace{17mm}
{\large Подготовка к ЕГЭ по профильной математике\par}
\vspace{16mm}
{\Large Анна Баннова\par}
\vspace{12mm}
{\large 1 ноября 2026 г.\par}
\vspace{10mm}
{Лёвин Артём Александрович\par}
{эксклюзивно для Анны Банновой\par}
\vspace{12mm}
\begin{tikzpicture}[x=1cm,y=1cm]
\draw[accent,line width=0.75pt] (0,0) .. controls (1.4,0.65) and (3.0,-0.4) .. (4.4,0)
.. controls (5.8,0.45) and (7.2,-0.48) .. (9,0.05);
\fill[darkaccent] (4.4,0) circle (1.4pt);
\end{tikzpicture}
\vspace{9mm}
\end{center}
\clearpage

\lessonsection{Повторение: подстановка в прикладную формулу}

\key{Главная идея.} Сначала переведите текст условия на язык величин: что известно, какие показатели должны быть максимальными и какую переменную нужно найти. Только после этого выполняйте вычисления.

\example{Пример 1. Рейтинг публикации}

Рейтинг вычисляется по формуле
\[
R=\frac{3In+Op+2Tr}{A},\qquad A\ne0.
\]
Все три показателя лежат в промежутке от $-2$ до $2$. Найдите $A$, если показатели максимальны, а рейтинг равен $30$.

\textbf{Шаг 1.} Максимально возможное значение каждого показателя равно $2$:
\[
In=2,\qquad Op=2,\qquad Tr=2,\qquad R=30.
\]
\textbf{Шаг 2.} Подставим данные и приведём числитель:
\[
30=\frac{3\cdot2+2+2\cdot2}{A}=\frac{12}{A}.
\]
\textbf{Шаг 3.} Умножим обе части равенства на $A$ и разделим на $30$:
\[
30A=12,\qquad A=\frac{12}{30}=\boxed{0{,}4}.
\]
\textbf{Проверка:} $12:0{,}4=30$.

\key{Памятка.} «Максимальны» означает: возьмите верхнюю границу указанного промежутка. При подстановке сохраняйте коэффициенты $3$ и $2$; не забывайте, что знаменатель исходной формулы должен быть отличен от нуля.

\lessonsection{Три величины в задачах на движение}

\begin{center}
\renewcommand{\arraystretch}{1.25}
\begin{tabular}{@{}>{\raggedright\arraybackslash}p{39mm}>{\raggedright\arraybackslash}p{38mm}>{\raggedright\arraybackslash}p{53mm}@{}}
\toprule
Величина & Обозначение & Формула\\
\midrule
Расстояние & $s$, км & $s=vt$\\
Скорость & $v$, км/ч & $v=\dfrac{s}{t}$, $t>0$\\
Время & $t$, ч & $t=\dfrac{s}{v}$, $v>0$\\
\bottomrule
\end{tabular}
\end{center}

Если два автомобиля движутся навстречу друг другу и встретились, то сумма пройденных ими расстояний равна начальному расстоянию между городами:
\[
\boxed{s_1+s_2=S}.
\]

\begin{center}
\begin{tikzpicture}[x=0.93cm,y=0.93cm,>=Latex]
\coordinate (A) at (0,0); \coordinate (M) at (5,0); \coordinate (B) at (10,0);
\draw[black!60,line width=0.65pt] (A)--(B);
\fill[accent] (A) circle (2.3pt); \fill[accent] (M) circle (2.3pt); \fill[accent] (B) circle (2.3pt);
\node[below=3pt] at (A) {А}; \node[below=3pt] at (M) {Место встречи}; \node[below=3pt] at (B) {Б};
\draw[->,accent,thick] (0.55,0.36)--(4.25,0.36);
\draw[->,accent,thick] (9.45,0.36)--(5.75,0.36);
\node[above] at (2.45,0.4) {$s_1=v_1t_1$};
\node[above] at (7.55,0.4) {$s_2=v_2t_2$};
\draw[<->,black!50] (0,-0.58)--(10,-0.58);
\node[below] at (5,-0.58) {$S=s_1+s_2$};
\end{tikzpicture}
\end{center}

\textbf{Если выехали одновременно,} то до момента встречи они находились в пути одинаковое время: $t_1=t_2$.

\textbf{Если один выехал раньше,} то до встречи он находился в пути дольше. При разнице старта $\Delta t$:
\[
\boxed{t_{\text{раньше}}=t_{\text{позже}}+\Delta t}.
\]

\lessonsection{Как составить таблицу и уравнение}

\key{Рабочий порядок:} прочитайте задачу, при необходимости сделайте схематический рисунок, обозначьте неизвестную, заполните таблицу и составьте уравнение по ещё не использованному условию.

\example{Пример 2. Одновременный выезд; найти время}

Из городов А и Б, расстояние между которыми $560$ км, одновременно навстречу друг другу выехали два автомобиля со скоростями $60$ и $80$ км/ч. Через сколько часов они встретятся?

Пусть $x$ ч --- время каждого автомобиля. Тогда:
\[
\begin{array}{lccc}
\toprule
&v\ (\text{км/ч})&t\ (\text{ч})&s\ (\text{км})\\
\midrule
\text{Из А}&60&x&60x\\
\text{Из Б}&80&x&80x\\
\bottomrule
\end{array}
\]
\textbf{Составляем уравнение по расстоянию:}
\[
60x+80x=560,\qquad140x=560,\qquad\boxed{x=4}.
\]
\textbf{Ответ:} через $4$ часа.

\clearpage
\example{Пример 3. Найти расстояние, пройденное первым автомобилем}

Расстояние между городами $480$ км. Скорости встречных автомобилей $55$ и $65$ км/ч. Они выехали одновременно. На каком расстоянии от города А произойдёт встреча?

\textbf{Способ 1.} Пусть $x$ км --- искомое расстояние от А. Тогда второй проедет $480-x$ км. Так как времена одинаковы,
\[
\frac{x}{55}=\frac{480-x}{65}.
\]
Отсюда $65x=55(480-x)$, то есть $120x=26400$ и $\boxed{x=220}$ км.

\textbf{Способ 2.} Возьмём за $x$ время движения. Тогда
\[
55x+65x=480,\qquad x=4\text{ ч}.
\]
Из А проехали $55\cdot4=\boxed{220}$ км.

\key{Вывод.} Можно выбрать в качестве $x$ величину, которую удобнее найти. Если выбрали время, а спрашивают расстояние, после решения уравнения обязательно выполните последний шаг: $s=vt$.

\clearpage
\begin{center}
{\Large\color{darkaccent}\bfseries Алгоритм решения текстовой задачи}\par
\vspace{8mm}
\end{center}
\tikzset{
startstop/.style={ellipse,draw=accent,very thick,align=center,minimum width=52mm,minimum height=9mm,inner sep=3pt,font=\small},
process/.style={rectangle,rounded corners=2pt,draw=accent!85,very thick,align=center,text width=94mm,minimum height=11mm,inner sep=4pt,font=\small},
decide/.style={diamond,aspect=2.2,draw=accent!90,very thick,align=center,inner sep=1.2pt,text width=34mm,font=\small},
arr/.style={-{Latex[length=2.2mm]},thick,draw=darkaccent}
}
\begin{center}
\begin{tikzpicture}[node distance=7mm and 12mm]
\node[startstop] (start) {Прочитайте условие};
\node[process,below=of start] (draw) {При необходимости сделайте схему движения};
\node[process,below=of draw] (var) {Выберите $x$: искомую или удобную для расчёта величину};
\node[process,below=of var] (table) {Запишите таблицу: скорость, время, расстояние};
\node[process,below=of table] (formula) {Заполните оставшийся столбец по формулам $s=vt$, $t=s/v$};
\node[process,below=of formula] (equation) {Составьте уравнение из условия и решите его};
\node[decide,below=10mm of equation] (choose) {Найденный $x$ --- то, что спрашивают?};
\node[startstop,below=15mm of choose] (yes) {Проверьте и запишите ответ};
\node[process,text width=36mm,right=12mm of choose] (calc) {Найдите искомую величину, подставив $x$};
\node[startstop,below=13mm of calc] (finish) {Проверьте и запишите ответ};
\draw[arr] (start)--(draw); \draw[arr] (draw)--(var);
\draw[arr] (var)--(table); \draw[arr] (table)--(formula);
\draw[arr] (formula)--(equation); \draw[arr] (equation)--(choose);
\draw[arr] (choose.south)--node[right,font=\small]{да}(yes.north);
\draw[arr] (choose.east)--node[above,font=\small]{нет}(calc.west);
\draw[arr] (calc)--(finish);
\end{tikzpicture}
\end{center}
\clearpage

\lessonsection{Выезд в разное время}

\example{Пример 4. Найти время после выезда второго автомобиля}

Расстояние между городами $580$ км. Первый автомобиль выехал из А со скоростью $80$ км/ч, а через $2$ часа из Б ему навстречу выехал второй со скоростью $60$ км/ч. Через сколько часов после выезда второго автомобиля они встретятся?

Пусть $x$ ч --- время второго. Первый находился в дороге на $2$ часа дольше.
\[
\begin{array}{lccc}
\toprule
&v\ (\text{км/ч})&t\ (\text{ч})&s\ (\text{км})\\
\midrule
\text{Первый}&80&x+2&80(x+2)\\
\text{Второй}&60&x&60x\\
\bottomrule
\end{array}
\]
Так как автомобили встретились, пройденные расстояния в сумме равны $580$ км:
\[
80(x+2)+60x=580,\quad 140x+160=580,\quad \boxed{x=3}.
\]
\textbf{Проверка:} первый ехал $5$ часов и прошёл $400$ км; второй ехал $3$ часа и прошёл $180$ км. Сумма $580$ км.

\example{Пример 5. Если за $x$ взять время первого автомобиля}

Расстояние между городами $510$ км. Первый автомобиль выехал со скоростью $60$ км/ч; второй стартовал навстречу ему через $2$ часа со скоростью $70$ км/ч. Найдите расстояние от города Б до места встречи.

Пусть $x$ ч --- время первого. Тогда второй находился в пути $x-2$ ч. Поскольку второй успел выехать до встречи, $x>2$.
\[
60x+70(x-2)=510,
\]
\[
130x-140=510,\qquad x=5\text{ ч}.
\]
Второй автомобиль проехал $70(5-2)=\boxed{210}$ км. Именно это расстояние спрашивают в условии.

\key{Сравните два способа записи времени.}
\begin{center}
\begin{tabularx}{\linewidth}{@{}X X@{}}
\toprule
Если $x$ --- время \textbf{второго} & Если $x$ --- время \textbf{первого}\\
\midrule
Первый: $x+2$, второй: $x$. & Первый: $x$, второй: $x-2$.\\
Условие: $x>0$. & Условие: $x>2$.\\
\bottomrule
\end{tabularx}
\end{center}

\lessonsection{На что обращать внимание}
\begin{enumerate}[label=\textcolor{accent}{\arabic*.}]
\item \textbf{«Одновременно»} означает равенство времён движения до встречи.
\item \textbf{«На два часа раньше»} означает большее время движения на два часа. Знак зависит от выбранного $x$.
\item \textbf{«На каком расстоянии от города Б?»} означает путь, пройденный автомобилем, стартовавшим из Б.
\item \textbf{Последний столбец таблицы} заполняйте по формуле, затем составляйте уравнение по условию задачи.
\item \textbf{Проверка:} время и расстояния должны быть положительными, а сумма пройденных путей --- равняться расстоянию между городами.
\item \textbf{Сложные дроби в уравнении?} Попробуйте взять за $x$ время, а искомое расстояние вычислить в конце.
\end{enumerate}

\clearpage
\begin{center}
{\Large\color{darkaccent}\bfseries Тренировочный блок}\par
{\small Путь А \textbullet\ 10 заданий \textbullet\ от простого к сложному}\par
\end{center}
\vspace{2mm}
\begingroup\setstretch{1.25}\setlength{\parskip}{1pt}
\noindent\key{Указание.} В заданиях 3--10 сначала оформите таблицу «скорость --- время --- расстояние», затем составьте уравнение. Если автомобили выехали не одновременно, отдельно укажите, кто находился в пути дольше.

\task{1}{Рейтинг публикации задаётся формулой
$R=\dfrac{3In+Op+2Tr}{A}$, где все показатели лежат от $-2$ до $2$ включительно. При максимальных показателях рейтинг равен $24$. Найдите $A$.}

\task{2}{По той же формуле вычисляют рейтинг, но теперь каждый показатель лежит от $-3$ до $3$. При максимальных значениях всех трёх показателей $R=15$. Найдите $A$.}

\task{3}{Из городов А и Б, расстояние между которыми $420$ км, одновременно навстречу друг другу выехали автомобили со скоростями $60$ и $80$ км/ч. Через сколько часов они встретятся?}

\task{4}{Два автомобиля одновременно выехали навстречу друг другу из городов, находящихся на расстоянии $600$ км. Их скорости $70$ и $80$ км/ч. Через сколько часов произойдёт встреча?}

\task{5}{Между городами А и Б $560$ км. Из А со скоростью $50$ км/ч и из Б со скоростью $90$ км/ч одновременно выехали автомобили навстречу друг другу. На каком расстоянии от А они встретятся?}

\task{6}{Расстояние между городами $600$ км. Из А навстречу автомобилю из Б одновременно выехали два автомобиля со скоростями $80$ и $70$ км/ч соответственно. На каком расстоянии от Б они встретятся?}

\task{7}{Между городами $430$ км. Первый автомобиль выехал навстречу второму со скоростью $70$ км/ч на $1$ час раньше. Скорость второго $50$ км/ч. Через сколько часов после старта второго автомобиля произойдёт встреча?}

\task{8}{Расстояние между А и Б $540$ км. Из А выехал автомобиль со скоростью $60$ км/ч, а через $2$ часа из Б навстречу ему --- автомобиль со скоростью $80$ км/ч. На каком расстоянии от Б они встретятся?}

\task{9}{Города А и Б разделяют $645$ км. Из А выехал автомобиль со скоростью $55$ км/ч; через $3$ часа из Б навстречу ему выехал автомобиль со скоростью $65$ км/ч. На каком расстоянии от А они встретятся?}

\task{10}{Между городами $655$ км. Первый автомобиль выехал из А со скоростью $65$ км/ч, а через $2$ часа второй выехал навстречу из Б со скоростью $85$ км/ч. Найдите расстояние от Б до места встречи. Ответ запишите в километрах.}

\endgroup
\clearpage
\begin{center}
{\Large\color{darkaccent}\bfseries Ответы}\par
\vspace{6mm}
\end{center}
\renewcommand{\arraystretch}{1.45}
\begin{center}
\begin{tabularx}{0.84\linewidth}{@{}>{\centering\arraybackslash}p{18mm}X@{}}
\toprule
№ & Ответ\\
\midrule
\answerrow{1}{$0{,}5$}
\answerrow{2}{$1{,}2$}
\answerrow{3}{$3$ ч}
\answerrow{4}{$4$ ч}
\answerrow{5}{$200$ км от А}
\answerrow{6}{$280$ км от Б}
\answerrow{7}{$3$ ч после выезда второго}
\answerrow{8}{$240$ км от Б}
\answerrow{9}{$385$ км от А}
\answerrow{10}{$297{,}5$ км от Б}
\bottomrule
\end{tabularx}
\end{center}
\vspace{14mm}
\begin{center}
{\color{darkaccent}\large Короткая памятка}\par
\vspace{4mm}
\begin{tabularx}{0.95\linewidth}{@{}>{\raggedright\arraybackslash}p{56mm}X@{}}
\toprule
Сигнал в условии & Что записать\\
\midrule
Выехали одновременно & $t_1=t_2$\\
Один выехал на $\Delta t$ раньше & $t_{\text{раньше}}=t_{\text{позже}}+\Delta t$\\
Встретились, двигались навстречу & $s_1+s_2=S$\\
Известны скорость и время & $s=vt$\\
Известны путь и скорость & $t=s/v$\\
Просили путь, а нашли время & Найдите $s=vt$ и запишите единицы\\
\bottomrule
\end{tabularx}
\end{center}
\end{document}
